\documentclass{article}
\usepackage{amsmath}

\begin{document}
	
	\title{Longitud de Arco de una Curva en el Plano xy}
	\author{Por: David Duarte}
	\date{\today}
	
	\maketitle
	
	\section*{Introducción}
	La longitud de arco de una curva en el plano xy es un concepto fundamental en cálculo y geometría diferencial. Nos permite medir la "longitud" de una curva en el plano, lo que es relevante en muchas aplicaciones, como en física, ingeniería y modelado matemático.
	
	\section*{Definición de Curva Paramétrica en el Plano xy}
	Una curva paramétrica en el plano xy es una representación de una curva en términos de una variable paramétrica, comúnmente denotada como $t$. La ecuación paramétrica de una curva en el plano xy toma la forma $\mathbf{r}(t) = (x(t), y(t))$, donde $x(t)$ y $y(t)$ son funciones que describen la posición de la curva en términos de $t$.
	
	\section*{Cálculo de la Longitud de Arco}
	La fórmula para calcular la longitud de arco de una curva paramétrica en el plano xy es dada por:
	\[ L = \int_{a}^{b} \sqrt{(x'(t))^2 + (y'(t))^2} \, dt \]
	donde $a \leq t \leq b$ es el intervalo de la curva paramétrica. Esta integral representa la suma de las longitudes de los segmentos tangentes a la curva, lo que nos da la longitud total de la curva.
	
	\section*{Ejemplo de Cálculo de Longitud de Arco}
	Para ilustrar el cálculo de la longitud de arco, consideremos la curva paramétrica $\mathbf{r}(t) = (t^2, 2t)$ en el intervalo $-2 \leq t \leq 2$. Calcularemos la longitud de arco utilizando la fórmula de la integral y discutiremos su significado geométrico.
	
	\section*{Aplicaciones y Significado Geométrico}
	La longitud de arco de una curva en el plano xy tiene aplicaciones en diversas áreas, como en la estimación de distancias recorridas, la determinación de la "curvatura" de una curva y en el modelado de trayectorias en el espacio bidimensional.
	

	
\end{document}
